\documentclass[12pt,a4paper]{article}

\usepackage[margin=1in]{geometry}
\usepackage{amsmath,amssymb}
\usepackage{xcolor}
\usepackage{enumitem}
\usepackage{hyperref}
\usepackage{array}
\usepackage{booktabs}
\usepackage{listings}

\hypersetup{
    colorlinks=true,
    linkcolor=blue,
    urlcolor=blue
}

\lstset{
    basicstyle=\ttfamily\small,
    breaklines=true,
    frame=single,
    columns=fullflexible
}

\title{\textbf{Machine Coding}\\
\large Course Tracker \& Exam Notes}
\author{}
\date{}

\begin{document}

\maketitle
\tableofcontents
\newpage

% =========================================================
\section{Machine Coding Overview}

\subsection{What is Machine Coding?}

Machine coding is the process of designing and implementing a working software system from a set of requirements, usually within a limited amount of time.

The goal is not only to make the program work, but to write code that is:

\begin{itemize}
    \item Correct
    \item Readable
    \item Maintainable
    \item Extensible
    \item Modular
    \item Testable
\end{itemize}

A machine coding problem typically gives you a problem statement and expects you to design classes, relationships, interfaces, and working functionality.

\subsection{Exam Definition}

\textbf{Machine coding:} A practical software-design exercise where requirements are converted into a working implementation while maintaining good object-oriented design and code quality.

% =========================================================
\section{Machine Coding vs Normal Coding}

Normal programming questions often focus mainly on:

\begin{itemize}
    \item Algorithm
    \item Data structures
    \item Time complexity
    \item Output correctness
\end{itemize}

Machine coding focuses more heavily on:

\begin{itemize}
    \item Class design
    \item Object-oriented principles
    \item Separation of responsibilities
    \item Extensibility
    \item Design patterns
    \item Interfaces
    \item Relationships between objects
\end{itemize}

\[
\boxed{\text{Machine Coding = Working Code + Good Design}}
\]

% =========================================================
\section{Typical Machine Coding Workflow}

A useful workflow is:

\[
Requirements
\rightarrow
Identify Objects
\rightarrow
Define Responsibilities
\rightarrow
Design Classes
\rightarrow
Define Relationships
\rightarrow
Implement
\rightarrow
Test
\]

\subsection{Step 1: Understand Requirements}

Read the problem carefully.

Extract:

\begin{itemize}
    \item What the system should do
    \item What entities exist
    \item What operations are required
    \item What constraints exist
    \item What can change in the future
\end{itemize}

Do not start writing code immediately.

\subsection{Step 2: Identify Objects}

Look for important nouns in the requirements.

For example, in a parking lot:

\begin{verbatim}
ParkingLot
Vehicle
ParkingSpot
Ticket
Payment
\end{verbatim}

These are possible classes.

\subsection{Step 3: Identify Responsibilities}

For every class ask:

\begin{quote}
What should this class be responsible for?
\end{quote}

Example:

\begin{verbatim}
Vehicle
    -> stores vehicle information

ParkingSpot
    -> stores spot information
    -> knows whether it is occupied

ParkingLot
    -> manages parking spots
    -> assigns spots
\end{verbatim}

\subsection{Step 4: Define Relationships}

Determine how objects interact.

Examples:

\begin{verbatim}
ParkingLot HAS ParkingSpot
ParkingSpot OCCUPIED BY Vehicle
ParkingLot CREATES Ticket
\end{verbatim}

\subsection{Step 5: Define Interfaces}

If multiple classes can perform the same type of operation, consider an interface or abstract class.

Example:

\begin{verbatim}
Payment
    |
    +-- CashPayment
    +-- CardPayment
    +-- UPIPayment
\end{verbatim}

\subsection{Step 6: Implement}

Implement the classes and their interactions.

\subsection{Step 7: Test}

Test:

\begin{itemize}
    \item Normal cases
    \item Edge cases
    \item Invalid inputs
    \item Multiple objects
    \item State changes
\end{itemize}

% =========================================================
\section{Object-Oriented Thinking in Machine Coding}

Machine coding strongly depends on object-oriented design.

\subsection{Encapsulation}

Keep data and the operations that manipulate it together.

Example:

\begin{lstlisting}[language=C++]
class BankAccount {
private:
    double balance;

public:
    void deposit(double amount) {
        if (amount > 0)
            balance += amount;
    }

    double getBalance() {
        return balance;
    }
};
\end{lstlisting}

The balance cannot be modified directly from outside the class.

\subsection{Abstraction}

Expose what an object does rather than how it does it.

Example:

\begin{lstlisting}[language=C++]
class Payment {
public:
    virtual void pay(double amount) = 0;
};
\end{lstlisting}

The caller only needs to know that payment can be performed.

\subsection{Polymorphism}

Different classes can implement the same interface differently.

\begin{lstlisting}[language=C++]
class CardPayment : public Payment {
public:
    void pay(double amount) override {
        cout << "Card payment";
    }
};

class UPIPayment : public Payment {
public:
    void pay(double amount) override {
        cout << "UPI payment";
    }
};
\end{lstlisting}

% =========================================================
\section{Class Responsibility}

One of the most important machine coding skills is deciding what belongs inside each class.

A class should have a clear responsibility.

Bad design:

\begin{verbatim}
ParkingLot
    -> stores parking data
    -> calculates payment
    -> sends notifications
    -> processes user authentication
    -> generates reports
\end{verbatim}

This creates a large class with many responsibilities.

Better:

\begin{verbatim}
ParkingLot
PaymentService
NotificationService
AuthenticationService
ReportService
\end{verbatim}

Each component has a clearer responsibility.

This connects directly to the \textbf{Single Responsibility Principle (SRP)}.

% =========================================================
\section{Composition vs Inheritance}

Machine coding often requires deciding between inheritance and composition.

\subsection{Inheritance}

Use inheritance when there is a genuine:

\[
\boxed{\text{IS-A}}
\]

relationship.

Example:

\begin{verbatim}
Vehicle
   |
   +-- Car
   +-- Bike
\end{verbatim}

A Car is a Vehicle.

\subsection{Composition}

Use composition when there is a:

\[
\boxed{\text{HAS-A}}
\]

relationship.

Example:

\begin{verbatim}
Car
 |
 +-- Engine
\end{verbatim}

A Car has an Engine.

\subsection{Rule of Thumb}

\[
IS-A \rightarrow Inheritance
\]

\[
HAS-A \rightarrow Composition
\]

Composition is often preferred when behavior needs to change independently.

% =========================================================
\section{Interfaces and Abstract Classes}

When multiple implementations should follow the same contract, use an interface-like abstract class.

Example:

\begin{lstlisting}[language=C++]
class Notification {
public:
    virtual void send(string message) = 0;
    virtual ~Notification() {}
};
\end{lstlisting}

Implementations:

\begin{lstlisting}[language=C++]
class EmailNotification : public Notification {
public:
    void send(string message) override {
        cout << "Email: " << message;
    }
};

class SMSNotification : public Notification {
public:
    void send(string message) override {
        cout << "SMS: " << message;
    }
};
\end{lstlisting}

Now another class can depend on the abstraction rather than a concrete implementation.

% =========================================================
\section{SOLID in Machine Coding}

SOLID principles help keep machine coding solutions maintainable.

\subsection{S -- Single Responsibility Principle}

A class should have one primary responsibility.

\[
\boxed{\text{One class should not handle unrelated responsibilities}}
\]

Example:

Do not make one class responsible for:

\begin{verbatim}
Payment + Notification + Logging + Authentication
\end{verbatim}

Separate them where appropriate.

\subsection{O -- Open/Closed Principle}

Software should be open for extension but closed for unnecessary modification.

Example:

If you support:

\begin{verbatim}
CardPayment
UPIPayment
\end{verbatim}

adding:

\begin{verbatim}
CashPayment
\end{verbatim}

should ideally not require changing the core checkout logic.

\subsection{L -- Liskov Substitution Principle}

A subclass should be usable wherever its base class is expected without breaking the expected behavior.

\subsection{I -- Interface Segregation Principle}

Prefer small, focused interfaces instead of one large interface containing unrelated methods.

\subsection{D -- Dependency Inversion Principle}

High-level classes should depend on abstractions rather than concrete implementations.

\[
HighLevel \rightarrow Interface \leftarrow LowLevel
\]

% =========================================================
\section{Design Patterns in Machine Coding}

Design patterns can provide reusable structures for common design problems.

Common examples:

\begin{center}
\begin{tabular}{p{0.25\textwidth} p{0.60\textwidth}}
\toprule
\textbf{Pattern} & \textbf{Typical Use}\\
\midrule
Factory & Creating different implementations\\
Strategy & Switching algorithms/behaviors\\
Observer & Notifying multiple objects\\
Builder & Constructing complex objects\\
Singleton & Ensuring a single shared instance\\
Adapter & Connecting incompatible interfaces\\
Decorator & Adding behavior dynamically\\
Command & Encapsulating actions\\
\bottomrule
\end{tabular}
\end{center}

Do not force a design pattern into a problem just because you know it. Use a pattern when it solves an actual design problem.

% =========================================================
\section{Example: Parking Lot}

Parking Lot is a common machine coding problem.

\subsection{Possible Classes}

\begin{verbatim}
ParkingLot
ParkingFloor
ParkingSpot
Vehicle
Ticket
Payment
\end{verbatim}

Possible vehicle hierarchy:

\begin{verbatim}
Vehicle
   |
   +-- Car
   +-- Bike
   +-- Truck
\end{verbatim}

Possible relationships:

\begin{verbatim}
ParkingLot
    |
    +-- ParkingFloor
            |
            +-- ParkingSpot

ParkingSpot
    |
    +-- Vehicle
\end{verbatim}

\subsection{Important Design Questions}

Before coding, decide:

\begin{itemize}
\item What types of vehicles are supported?
\item What types of parking spots exist?
\item How is a spot selected?
\item How is a ticket generated?
\item How is parking duration calculated?
\item How is payment calculated?
\item What happens when the parking lot is full?
\end{itemize}

The goal is to turn these requirements into clear responsibilities.

% =========================================================
\section{Example: Tic Tac Toe}

Possible classes:

\begin{verbatim}
Game
Board
Player
Symbol
\end{verbatim}

Responsibilities:

\begin{verbatim}
Board
    -> stores cells
    -> places symbols

Player
    -> stores player information
    -> chooses a move

Game
    -> controls game flow
    -> checks game state
\end{verbatim}

Avoid putting everything inside one \texttt{Game} class.

% =========================================================
\section{Example: LRU Cache}

An LRU cache needs to support operations such as:

\begin{verbatim}
get(key)
put(key, value)
\end{verbatim}

with efficient access and removal of the least recently used item.

A common implementation uses:

\begin{itemize}
\item Hash map for fast lookup
\item Doubly linked list for tracking usage order
\end{itemize}

Conceptually:

\[
HashMap + DoublyLinkedList
\]

The machine coding focus is not only the data structure but also how the responsibilities are organized into classes.

% =========================================================
\section{Extensibility}

A good machine coding solution should make reasonable future changes easy.

Suppose a payment system supports:

\begin{verbatim}
UPI
Card
\end{verbatim}

Later the system needs:

\begin{verbatim}
Cash
Wallet
NetBanking
\end{verbatim}

A design based on an abstraction:

\begin{verbatim}
PaymentStrategy
    |
    +-- UPIPayment
    +-- CardPayment
    +-- CashPayment
    +-- WalletPayment
\end{verbatim}

can be extended without rewriting the whole system.

% =========================================================
\section{Dependency Injection}

Dependency Injection means providing an object with the dependencies it needs rather than making the object create them internally.

Without dependency injection:

\begin{lstlisting}[language=C++]
class Checkout {
private:
    CardPayment payment;
};
\end{lstlisting}

Now Checkout is tightly coupled to CardPayment.

With dependency injection:

\begin{lstlisting}[language=C++]
class Checkout {
private:
    Payment* payment;

public:
    Checkout(Payment* p) : payment(p) {}
};
\end{lstlisting}

Now:

\begin{verbatim}
Checkout
   |
   v
Payment interface
   ^
   |
CardPayment / UPIPayment
\end{verbatim}

This improves flexibility and testability.

% =========================================================
\section{Testing Machine Coding Solutions}

After implementation, test the behavior.

\subsection{Functional Tests}

Check normal operations.

Example:

\begin{verbatim}
create vehicle
park vehicle
generate ticket
exit vehicle
calculate payment
\end{verbatim}

\subsection{Edge Cases}

Examples:

\begin{itemize}
\item Parking lot is full.
\item Invalid vehicle type.
\item Duplicate ID.
\item Empty input.
\item Object does not exist.
\item Invalid state transition.
\end{itemize}

\subsection{State Testing}

For state-based systems, check transitions.

Example:

\[
Idle \rightarrow Running \rightarrow Completed
\]

Verify that invalid transitions are rejected when required.

% =========================================================
\section{Common Mistakes}

\begin{enumerate}
\item Starting to code before understanding requirements.
\item Creating one giant class.
\item Using inheritance everywhere.
\item Making every variable public.
\item Hard-coding behavior that should be extensible.
\item Mixing unrelated responsibilities.
\item Overusing design patterns.
\item Ignoring invalid inputs and edge cases.
\item Not testing interactions between classes.
\item Writing code that works only for the initial example.
\end{enumerate}

% =========================================================
\section{Machine Coding Checklist}

Before coding:

\begin{itemize}[label=$\square$]
\item Understand all requirements
\item Identify entities/objects
\item Identify responsibilities
\item Identify relationships
\item Decide inheritance vs composition
\item Identify possible abstractions
\item Identify behavior that may change
\item Choose patterns only if useful
\end{itemize}

While coding:

\begin{itemize}[label=$\square$]
\item Keep classes focused
\item Encapsulate data
\item Program to abstractions
\item Avoid unnecessary coupling
\item Use meaningful names
\item Keep methods focused
\item Handle invalid states
\end{itemize}

After coding:

\begin{itemize}[label=$\square$]
\item Test normal cases
\item Test edge cases
\item Test invalid inputs
\item Check extensibility
\item Review SOLID principles
\item Remove unnecessary complexity
\end{itemize}

% =========================================================
\section{Exam-Oriented Questions}

\subsection*{What is machine coding?}
A practical software-design exercise where requirements are converted into a working implementation while maintaining good object-oriented design.

\subsection*{Why is OOP important in machine coding?}
OOP allows the system to be divided into objects with clear responsibilities, making the design easier to understand, maintain, test, and extend.

\subsection*{When should inheritance be used?}
When there is a genuine ``IS-A'' relationship and the subclass follows the abstraction represented by the base class.

\subsection*{When should composition be used?}
When one object contains or uses another object, representing a ``HAS-A'' relationship.

\subsection*{Why are interfaces useful?}
They define contracts and allow multiple implementations while reducing coupling to concrete classes.

\subsection*{Why are SOLID principles important?}
They help create maintainable, flexible, testable, and extensible designs.

% =========================================================
\section{One-Page Memory Sheet}

\begin{center}
\textbf{\Large MACHINE CODING = REQUIREMENTS + DESIGN + IMPLEMENTATION}
\end{center}

\begin{verbatim}
1. Understand Requirements
          |
          v
2. Identify Objects
          |
          v
3. Assign Responsibilities
          |
          v
4. Define Relationships
          |
          v
5. Choose Abstractions
          |
          v
6. Implement
          |
          v
7. Test
\end{verbatim}

\textbf{Core rules:}

\begin{verbatim}
IS-A       -> Inheritance
HAS-A      -> Composition

One responsibility       -> SRP
Extend without rewriting -> OCP
Subtypes remain valid    -> LSP
Small interfaces         -> ISP
Depend on abstractions   -> DIP

Factory    -> Object creation
Strategy   -> Interchangeable behavior
Observer   -> Notifications
Builder    -> Complex construction
Adapter    -> Compatible interfaces
Decorator  -> Added behavior
Command    -> Encapsulated action
\end{verbatim}

\end{document}
